\chapter*{Abstract}
Battery degradation is governed by electrochemical processes evolving across multiple time scales and operating conditions. Electrochemical impedance spectroscopy (EIS) is widely used to study these processes because it provides frequency-resolved information on interfacial kinetics, transport phenomena, and charge storage mechanisms. However, classical EIS relies on stationary conditions and long acquisition times, which restrict its use during realistic battery operation and limit its applicability in long-term aging studies.

This work addresses these limitations by developing a framework for the estimation of instantaneous impedance under non-stationary conditions, referred to as dynamic electrochemical impedance spectroscopy (DEIS). The method is grounded in a system-theoretical interpretation of impedance as the linearized response of a time-varying electrochemical system, and relies on broadband multi-frequency perturbations and Fourier-based analysis to track impedance continuously during operation.

A dedicated experimental and software framework was developed to enable the design, execution, and interpretation of dynamic impedance experiments. The methodology introduces automated experiment management and online impedance estimation, enabling long-duration measurements and substantially reducing data storage requirements. The frequency bandwidth achieved for battery measurements spans up to seven decades, extending the range typically explored in dynamic impedance studies.

The applicability of DEIS is demonstrated across different electrochemical systems, chemistries, and cell configurations, including full cells, half-cells, and three-electrode setups. The results show that time-resolved impedance reveals the evolution of interfacial kinetics, transport processes, adsorption phenomena, and metal plating signatures during operation. In battery systems, these measurements highlight clear relationships between impedance features, state of charge, and degradation mechanisms, supporting the use of DEIS as a diagnostic tool for tracking state of health during aging.

\newpage
\chapter*{Zusammenfassung}
Die Alterung von Batterien wird durch elektrochemische Prozesse bestimmt, die sich über mehrere Zeitskalen und Betriebsbedingungen hinweg entwickeln. Die elektrochemische Impedanzspektroskopie (EIS) ist ein etabliertes Werkzeug zur Untersuchung dieser Prozesse, da sie frequenzaufgelöste Informationen über Grenzflächenkinetik, Transportphänomene und Ladungsspeichermechanismen liefert. Klassische EIS setzt jedoch stationäre Bedingungen und lange Messzeiten voraus, was ihren Einsatz unter realistischen Betriebsbedingungen einschränkt und ihre Anwendbarkeit in Langzeit-Alterungsstudien begrenzt.

Diese Arbeit adressiert diese Einschränkungen durch die Entwicklung eines Rahmens zur Bestimmung der momentanen Impedanz unter nichtstationären Bedingungen, bezeichnet als Dynamic Electrochemical Impedance Spectroscopy (DEIS). Die Methode basiert auf einer systemtheoretischen Interpretation der Impedanz als linearisierten Antwort eines zeitvarianten elektrochemischen Systems und nutzt breitbandige Mehrfrequenz-Anregungen sowie Fourier-basierte Auswerteverfahren, um die Impedanz kontinuierlich während des Betriebs zu verfolgen.

Zur Umsetzung wurde ein experimentelles und softwareseitiges Framework entwickelt, das die Planung, Durchführung und Interpretation dynamischer Impedanzexperimente ermöglicht. Die Methodik umfasst ein automatisiertes Experimentmanagement sowie eine Online-Schätzung der Impedanz, wodurch Langzeitmessungen realisierbar werden und der Speicherbedarf erheblich reduziert wird. Der für Batteriemessungen erreichte Frequenzbereich erstreckt sich über bis zu sieben Dekaden und erweitert damit den in dynamischen Impedanzstudien typischerweise untersuchten Bereich.

Die Anwendbarkeit von DEIS wird an verschiedenen elektrochemischen Systemen, Chemien und Zellkonfigurationen demonstriert, einschließlich Vollzellen, Halbzellen und Drei-Elektroden-Aufbauten. Die Ergebnisse zeigen, dass die zeitaufgelöste Impedanz die Entwicklung von Grenzflächenkinetik, Transportprozessen, Adsorptionsphänomenen und Metallabscheidung während des Betriebs sichtbar macht. In Batteriesystemen verdeutlichen diese Messungen klare Zusammenhänge zwischen Impedanzmerkmalen, Ladezustand und Alterungsmechanismen und unterstreichen damit das Potenzial von DEIS als diagnostisches Werkzeug zur Überwachung des Gesundheitszustands von Batterien im Betrieb.